% article-template.tex — Template for a research paper (amsart class)
%
% Compile: latexmk (uses .latexmkrc at the project root)
% Engine:  pdfLaTeX

\documentclass[12pt]{amsart}

% ═════════════════════════════════════════════════════════════════════════════
% ┃ ALL \usepackage, macro, and environment declarations live in                ┃
% ┃ article-template.sty (loaded below), which loads the vendored shared        ┃
% ┃ packages from packages/ and then this project's additions. Do NOT add       ┃
% ┃ \usepackage / \newcommand / \newtheorem here. This preamble carries ONLY    ┃
% ┃ the document class, that one style package, and semantic metadata           ┃
% ┃ (title / author / date / subjclass / keywords). See STYLE.md.               ┃
% ═════════════════════════════════════════════════════════════════════════════
\usepackage{article-template}

% ─── Document Metadata ──────────────────────────────────────────────────────
\title{Title of the Paper}

\author{Author One}
\address{Department of Mathematics, University of Example, City, State, ZIP}
\email{author.one@example.edu}

\author{Author Two}
\address{Department of Mathematics, Another University, City, State, ZIP}
\email{author.two@example.edu}

\date{\today}

\subjclass[2020]{Primary XXXXX; Secondary YYYYY}
\keywords{keyword one, keyword two, keyword three}

% ═══════════════════════════════════════════════════════════════════════════════

%#
%# DOCUMENT: Title of the Paper
%#
\begin{document}

\begin{abstract}
  This paper investigates the properties of widgets
  and establishes a characterization of compact widgets
  in terms of sequential compactness.
  We prove that every compact widget is sequentially compact
  and discuss the converse under additional hypotheses.
\end{abstract}

\maketitle

% ─── Introduction ────────────────────────────────────────────────────────────

%#
%## SECTION: Introduction
%#
\section{Introduction}\label{sec:introduction}

The study of widgets was initiated by Doe and Smith~\cite{placeholder2025}.
In this paper, we establish the following result.

\begin{theorem}\label{thm:main}
  Every compact widget is sequentially compact.
\end{theorem}

The proof of \cref{thm:main} is given in \cref{sec:proof}.

%#
%### SUBSECTION: Organization
%#
\subsection*{Organization}

In \cref{sec:preliminaries}, we recall the necessary definitions.
\Cref{sec:proof} contains the proof of \cref{thm:main}.
We conclude in \cref{sec:remarks} with open questions.

%#
%### SUBSECTION: Acknowledgements
%#
\subsection*{Acknowledgements}

\todo{Add acknowledgements.}

% ─── Preliminaries ──────────────────────────────────────────────────────────

%#
%## SECTION: Preliminaries
%#
\section{Preliminaries}\label{sec:preliminaries}

We fix notation and recall the basic definitions.

\begin{definition}\label{def:widget}
  A \emph{widget} is a pair $(X, \tau)$ where $X$ is a set and
  $\tau \subseteq \mathcal{P}(X)$ satisfies:
  \begin{enumerate}
  \item
    $\emptyset, X \in \tau$;
  \item
    $\tau$ is closed under arbitrary unions;
  \item
    $\tau$ is closed under finite intersections.
  \end{enumerate}
\end{definition}

\begin{definition}\label{def:compact}
  A widget $(X, \tau)$ is \emph{compact} if every open cover of $X$ has a
  finite subcover.
\end{definition}

% ─── Main Proof ──────────────────────────────────────────────────────────────

%#
%## SECTION: Proof of the Main Theorem
%#
\section{Proof of the Main Theorem}\label{sec:proof}

The key ingredient is the following lemma.

\begin{lemma}\label{lem:subsequence}
  Let $(X, \tau)$ be a compact widget and let $(x_n)$ be a sequence in $X$.
  Then $(x_n)$ has a convergent subsequence.
\end{lemma}

\begin{proof}
  Suppose for contradiction that no subsequence converges.
  For each $x \in X$, choose an open neighborhood $U_x$ containing
  only finitely many terms of the sequence.
  Then $\{U_x\}_{x \in X}$ is an open cover with no finite subcover,
  contradicting compactness.
\end{proof}

\begin{proof}[Proof of \cref{thm:main}]
  This follows immediately from \cref{lem:subsequence}.
\end{proof}

% ─── Concluding Remarks ─────────────────────────────────────────────────────

%#
%## SECTION: Concluding Remarks
%#
\section{Concluding Remarks}\label{sec:remarks}

\begin{conjecture}\label{conj:converse}
  Every sequentially compact widget is compact.
\end{conjecture}

\Cref{conj:converse} is known to hold under additional set-theoretic
hypotheses;
see~\cite{example2024}.

% ─── Bibliography ────────────────────────────────────────────────────────────
\bibliographystyle{alpha}
\bibliography{references}

\end{document}

%%% Local Variables:
%%% mode: LaTeX
%%% TeX-master: t
%%% End:
